% Pista pro-25s-q3ab · Probabilitat
% Procedència: PAU Catalunya 2025 · Convocatòria de setembre · Sèrie 3 · Exercici 3
\documentclass[11pt,a4paper]{article}
\usepackage{pista}
\pistainfo{Catalunya 2025}{Convocatòria de setembre}{Sèrie 3}

\begin{document}

\section*{\textcolor{blauPista}{Exercici 3}}

\noindent La lesió per sesamoïditis (inflamació de l'os sesamoide del peu) és relativament habitual entre la població que practica esports d'impacte (atletisme, bàsquet, tennis\ldots). En una població d'esportistes, s'ha fet un estudi diferenciant entre els que practiquen esports d'impacte i els que practiquen esports sense impacte brusc (com ara natació, pilates, senderisme\ldots). S'ha pogut determinar que el $45\,\%$ practiquen esports d'impacte. Entre aquests, un $10\,\%$ pateixen lesions per sesamoïditis, mentre que entre els que no practiquen esports d'impacte només un $3\,\%$ presenten aquesta lesió. Escollim un esportista a l'atzar.

\subsection*{\textit{a)} \normalfont Quina és la probabilitat que pateixi sesamoïditis? \hfill\textnormal{[0,75~punts]}}

\pista{Un diagrama d'arbre t'ajudarà a visualitzar els dos camins que porten a patir sesamoïditis: o bé practica esports d'impacte i té la lesió, o bé no practica esports d'impacte i també la té. Aplica el teorema de la probabilitat total.}

\subsection*{\textit{b)} \normalfont Si l'esportista escollit té una lesió per sesamoïditis, quina és la probabilitat que practiqui esports d'impacte? \hfill\textnormal{[0,75~punts]}}

\pista{Aquest és un exemple clar on has d'aplicar el Teorema de Bayes, perquè et demanen una probabilitat condicionada en sentit \emph{invers} respecte de la informació que et donen: tu coneixes la probabilitat de ``tenir la lesió \emph{sabent que} practica esports d'impacte'', però et demanen la probabilitat de ``practicar esports d'impacte \emph{sabent que} té la lesió''.}

\end{document}
